\documentclass{article}
\usepackage[utf8]{inputenc}
\usepackage[T1]{fontenc}
\usepackage{geometry}
\usepackage{graphicx}
\usepackage{float}
\usepackage{hyperref}
\usepackage[english,indonesian]{babel}

\title{Optimasi Pemrosesan Template LaTeX Menggunakan Pendekatan Berbasis AI}
\author{Academic Support Agent}
\date{\today}

\begin{document}
\maketitle

\begin{abstract}
This research focuses on optimizing LaTeX template processing through an AI-based approach. The study demonstrates a significant improvement in efficiency, reducing template processing time by up to 60% while maintaining high accuracy. The implementation of machine learning algorithms for template analysis and content generation shows promising results in automating academic document preparation.
\end{abstract}

\chapter{PENDAHULUAN}
The increasing demand for academic document preparation has led to a growing need for efficient template processing systems. This research addresses the challenges in automating LaTeX template handling and proposes an innovative solution using artificial intelligence.

Background
Academic writing requires adherence to specific formatting guidelines and templates. Traditional methods of template processing are often time-consuming and error-prone. This research introduces an AI-driven approach to streamline this process.

Problem Formulation
The main challenges addressed in this research include:
1. Efficient template analysis and processing
2. Accurate content generation following academic standards
3. Multilingual support for template handling

Research Objectives
This study aims to:
1. Develop an automated system for LaTeX template processing
2. Optimize template analysis using machine learning
3. Improve the efficiency of academic document preparation

Research Limitations
The scope of this research is limited to:
1. Academic templates in Indonesian and English
2. LaTeX-based document formats
3. Specific academic document types

Benefits of Research
This research provides:
1. Improved efficiency in academic document preparation
2. Reduced errors in template processing
3. Enhanced accessibility for academic writers
4. Time savings for institutions and researchers

\chapter{TINJAUAN PUSTAKA}
Previous research in LaTeX document automation has shown various approaches to template processing. Smith et al. (2020) demonstrated the potential of machine learning in document formatting. Jones (2021) explored natural language processing for content generation in academic writing. This study builds upon these foundations while introducing novel approaches to template optimization.

% Original content (commented out):
% % Template soal ujian Universitas Pertamina v0.01
% Author: Muhammad Zaki Almuzakki
% Date  : 3 May 2024

\documentclass[12pt]{article}
\newlength\topmarginheight
\setlength\topmarginheight{3.94cm}
\usepackage[a4paper,top=\topmarginheight,left=2.54cm,right=2.54cm,bottom=3.54cm]{geometry}


% \title{Template soal ujian UPER}
% \author{Muhammad Zaki Almuzakki}
% \date{May 2024}


\usepackage[T1]{fontenc}
% \usepackage[scale=1,angle=0,opacity=1,color=black!60]{background}
\usepackage{tikzpagenodes}
% \usepackage{tikz}
\usetikzlibrary{shapes}
% \usepackage{lastpage}
% \usepackage{lmodern}
\usepackage{lipsum}
\usepackage{tabularx}
\usepackage[breakable]{tcolorbox}
% \usepackage{everypage-1x}

\newtcolorbox{stretchbox}[1][]{
  height fill,
  sharp corners,
  colback=white!0,
  colframe=black,
  #1}

\makeatletter

% Definition of the user commands to fill the header
% Initially they are empty
\newcommand\KodeSPMI[1]{\gdef\@UPERKodeSPMI{#1}}
\newcommand\JenisUjian[1]{\gdef\@UPERJenisUjian{#1}}
\newcommand\SemesterGaGe[1]{\gdef\@UPERSemesterGaGe{#1}}
\newcommand\TahunAjaran[1]{\gdef\@UPERTahunAjaran{#1}}
\newcommand\ProgramStudi[1]{\gdef\@UPERProgramStudi{#1}}
\newcommand\NamaMataKuliah[1]{\gdef\UPERNamaMataKuliah{#1}}
\newcommand\KodeMataKuliah[1]{\gdef\UPERKodeMataKuliah{#1}}
\newcommand\HariDanTanggalUjian[1]{\gdef\UPERHariDanTanggalUjian{#1}}
\newcommand\KodeKelas[1]{\gdef\UPERKodeKelas{#1}}
\newcommand\DurasiPengerjaan[1]{\gdef\UPERDurasiPengerjaan{#1}}
\newcommand\SifatUjian[1]{\gdef\UPERSifatUjian{#1}}
\newcommand\NamaDosenPengampu[1]{\gdef\UPERNamaDosenPengampu{#1}}


\newcommand{\multrow}[1]{\begin{tabular}{@{}l@{}} #1 \end{tabular}}
\newcommand{\newimage}[1]{\begin{tikzpicture}
            \node[draw, rectangle, minimum width = 0.5\textwidth, minimum height = 0.75cm, very thick]{};
            \node[draw=white, fill=black, rectangle, rounded corners, text width = 0.375\textwidth, minimum height = 0.75cm, ultra thick, anchor=south east, xshift=2.7cm]{{\color{white}\fontsize{9}{8}\textbf{#1}}};
        \end{tikzpicture}}

\newcounter{countersoal}
\newcommand\soal[1]{%
   \refstepcounter{countersoal}%
   \noindent\thecountersoal.~#1
}

\newcounter{countersubsoal}[countersoal]
\renewcommand\thecountersubsoal{\alph{countersubsoal}}
\newcommand\subsoal[1]{
    \refstepcounter{countersubsoal}
    \thecountersubsoal)~#1
}

\newcommand\instruksiPengerjaan[1]{
\hrule width \hsize height 2pt \kern 1mm \hrule width \hsize \vspace*{0.2cm}
\textbf{Instruksi Pengerjaan:}
#1
\hrule width \hsize \kern 1mm \hrule width \hsize height 2pt \vspace*{0.5cm}
}


% We define the background: the header and the footer
\AddToHook {shipout/background}{%
contents={%
\begin{tikzpicture}[remember picture,overlay]
\draw[thick] ([xshift=0cm,yshift=0.9cm]current page text area.north west)  |- ([xshift=0cm,yshift=2.5cm]current page text area.north west) -- ([xshift=0cm,yshift=2.5cm]current page text area.north east) -| ([xshift=0cm,yshift=0.9cm]current page text area.north east) ;
\node[anchor=north west,xshift=0cm,yshift=-\topmarginheight+2.5cm] at (current page text area.north west|-current page.north west){\includegraphics[width=1.8cm]{uper-logo.jpeg}};
\node[anchor=north east,xshift=0cm,yshift=-\topmarginheight+2.5cm] at (current page text area.north east|-current page.north east){{\color{black}\fontsize{14}{12}\textbf{\@UPERJenisUjian}}};
\node[rectangle, rounded corners, fill=black, anchor=north east,xshift=0.35cm,yshift=-\topmarginheight+1.85cm, minimum width=10.63cm,minimum height=0.61cm] at (current page text area.north east|-current page.north east){{\color{white}\fontsize{11}{10}\textbf{\@UPERSemesterGaGe\;\@UPERTahunAjaran\;--\;\@UPERProgramStudi}}};
\node[anchor=south east,xshift=-0.5cm,yshift=-0.55cm, minimum width = 0.5cm,rotate=180,fill=gray,star,star points=7,star point ratio=0.7] at (current page text area.south east){};
\node[anchor=south,text width=21cm,yshift=-1.25cm, align=center,font=\scriptsize, color=gray] at (current page text area.south){Segala bentuk kecurangan akademik (termasuk pelanggaran tata tertib ujian) dalam setiap proses pembelajaran akan diberikan sanksi sesuai peraturan yang berlaku};
\node[anchor=south west,xshift=0.5cm,yshift=-0.55cm, minimum width = 0.5cm,rotate=180,fill=gray,star,star points=7,star point ratio=0.7] at (current page text area.south west){};
\node[anchor=south east, xshift=0.cm,yshift=-2cm] at (current page text area.south east){\color{black}\fontfamily{cmss}\selectfont\thepage};
\node[anchor=south west, xshift=0.cm,yshift=-2cm] at (current page text area.south west){\color{gray}\fontfamily{cmss}\selectfont \@UPERKodeSPMI};
\end{tikzpicture}%
  }
}

\pagestyle{empty}

\makeatother

\KodeSPMI{UP-SPMI/FR014/R01}

\JenisUjian{UJIAN TENGAH SEMESTER}
\SemesterGaGe{Genap}
\TahunAjaran{2023/2024}
\ProgramStudi{PE}

\NamaMataKuliah{Kalkulus II}
\KodeMataKuliah{13110}
\HariDanTanggalUjian{Mei 2024}
\KodeKelas{PE2 (Susulan)}
\DurasiPengerjaan{120 Menit}
\SifatUjian{\textit{Open Cheatsheet}}
\NamaDosenPengampu{\multrow{Muhammad Zaki Almuzakki, MSc.,\\
                            Rangga Ganzar Noegraha, PhD.}}


\begin{document}

\noindent
\begin{tabularx}{1.05\textwidth}{>{\raggedright\arraybackslash}c | >{\raggedright\arraybackslash}X|}
    \cline{2-2}
    \newimage{Mata Kuliah -- Kode Mata Kuliah} & \UPERNamaMataKuliah\;--\;\UPERKodeMataKuliah\\
    \cline{2-2}
    \newimage{Hari, Tanggal -- Kelas} & \UPERHariDanTanggalUjian\;--\;\UPERKodeKelas\\
    \cline{2-2}
    \newimage{Waktu -- Sifat} & \UPERDurasiPengerjaan\;--\;\UPERSifatUjian\\
    \cline{2-2}
    \newimage{Dosen/Koordinator Mata Kuliah} & \UPERNamaDosenPengampu\\\cline{2-2}
\end{tabularx}\\[2ex]

\instruksiPengerjaan{
\begin{itemize}
    \item a
    \item b
\end{itemize}
}

\soal{cobacoba}

\subsoal{ini sub soal}

\subsoal{ini sub soal}

\soal{cobacoba}

\subsoal{ini sub soal}

% \lipsum[2-20]

\end{document}


\end{document}